\documentclass[11pt]{article}
\usepackage[a4paper,margin=1in]{geometry}
\usepackage[T1]{fontenc}
\usepackage{lmodern,microtype}
\usepackage{amsmath,amssymb,amsthm,mathtools}
\usepackage{booktabs,enumitem}
\usepackage{xcolor}
\usepackage[numbers,sort&compress]{natbib}
\usepackage[colorlinks=true,linkcolor=blue!55!black,citecolor=blue!55!black]{hyperref}

\newtheorem{theorem}{Theorem}[section]
\newtheorem{corollary}[theorem]{Corollary}
\theoremstyle{remark}
\newtheorem{remark}[theorem]{Remark}

\title{Finite-Jet Contraction and the All-Order Obstruction}
\author{Liang Wang}
\date{July 2026}

\begin{document}
\maketitle

\begin{abstract}
The trace-adaptive selector of RH-238 enforces a noise-dependent finite-jet
tolerance \citep{WangRH238}.  We prove the exact consequence.  If
$J_{m,R}(\tau_\sigma)\le\varepsilon_\sigma$ and
$\varepsilon_\sigma\to0$, then the residual jets are Cauchy in the
$J_{m,R}$ seminorm; in fact
\[
 J_{m,R}(\tau_\sigma-\tau_{\sigma'})
 \le\varepsilon_\sigma+\varepsilon_{\sigma'}.
\]

All 30 adjacent-scale and 16 dual-channel archived cases satisfy the
corresponding triangle bound.  The minimum slacks are $0.00246$ and $0.00220$.
The actual adjacent distances are not monotonically contracting, and between
one and ten shell prefixes can meet a given tolerance.  Most importantly,
fixed-$m$ Cauchy convergence does not control coefficients above order $m$.
Thus adaptive finite jets advance the coefficient route but do not imply a
locally uniform regularized determinant limit.
\end{abstract}

\section{Finite-jet topology}

For trace vectors through order $m$, define
\[
 J_{m,R}(\tau)=\sum_{n=2}^m\frac{|\tau_n|}{n}R^n.
\]
This is a seminorm because the first coordinate is ignored.  It is the
natural coefficient majorant for the truncated $\det_2$ logarithm.

\begin{theorem}[Tolerance contraction]\label{thm:contraction}
Suppose a family $\tau_\alpha$ satisfies
$J_{m,R}(\tau_\alpha)\le\varepsilon_\alpha$.  Then
\[
 J_{m,R}(\tau_\alpha-\tau_\beta)
 \le\varepsilon_\alpha+\varepsilon_\beta.
\]
If $\varepsilon_\alpha\to0$, the family converges to zero in the fixed
finite-jet seminorm.
\end{theorem}

\begin{proof}
The first inequality is the triangle inequality.  Taking $\alpha,\beta$ in a
tail where both tolerances are small proves the Cauchy property.  Taking
the assumed one-point bounds directly gives convergence to the zero jet.
\end{proof}

\begin{remark}
The zero limit reflects the normalization chosen in RH-238: the adaptive
cloud is selected to make the residual finite jet small.  It is not yet the
coefficient anchor required to identify a nontrivial deterministic numerator.
\end{remark}

\begin{corollary}[Coefficientwise finite-order control]\label{cor:coeff}
If $R>0$ and $J_{m,R}(\tau_\alpha)\le\varepsilon_\alpha$, then for every
$2\le n\le m$,
\[
 |\tau_{\alpha,n}|\le \frac{n\varepsilon_\alpha}{R^n},
 \qquad
 |\tau_{\alpha,n}-\tau_{\beta,n}|
 \le \frac{n(\varepsilon_\alpha+\varepsilon_\beta)}{R^n}.
\]
\end{corollary}

\begin{proof}
Each nonnegative weighted coefficient is bounded by the full defining sum;
apply the same observation to Theorem~\ref{thm:contraction}.
\end{proof}

Thus vanishing tolerances give simultaneous convergence of every coefficient
in the fixed finite window, with explicit rates inherited from the tolerance
schedule.

\section{Archived bounds}

For adjacent scales the theorem gives
$J(\tau_{\sigma_i}-\tau_{\sigma_{i+1}})\le\sigma_i+\sigma_{i+1}$.  For two
channels at the same scale it gives the bound $2\sigma$.

\begin{table}[ht]
\centering
\small
\begin{tabular}{lr}
\toprule
Quantity & value\\
\midrule
Adjacent cases & 30\\
Channel cases & 16\\
Minimum adjacent bound slack & $0.0024624$\\
Minimum channel bound slack & $0.0022040$\\
Maximum actual adjacent distance & $0.0231244$\\
Both last-four actual sequences strictly contract & no\\
Admissible-prefix count range & $1$--$10$\\
\bottomrule
\end{tabular}
\caption{Finite-jet contraction audit.}
\end{table}

The deterministic upper bounds decrease with the noise schedule, but the
actual distances fluctuate because the selected rank and shell phases jump.
This distinction matters: a theorem can provide Cauchy control even when a
short finite sequence is not pointwise monotone.

The last four left-channel distances happen to contract strictly, while the
right-channel sequence does not.  Neither observation changes the theorem:
monotonicity is not required for Cauchy convergence, and a short monotone tail
would not prove the all-order determinant statement.

\section{Why finite jets are insufficient}

For each fixed $m$, convergence of the first $m$ coefficients gives only
convergence of a Taylor polynomial.  A sequence of entire functions can have
all first $m$ coefficients equal and still grow arbitrarily at order $m+1$.
For example,
\[
 f_k(z)=\exp(kz^{m+1})
\]
has the same logarithmic jet through order $m$ for every $k$, but is not
locally bounded on any disk containing a nonzero point with positive real
$z^{m+1}$.

Therefore the implication
\[
 \text{vanishing fixed finite jet}
 \Longrightarrow
 \text{normal entire family}
\]
is false.  Uniformity in the order is indispensable.

\begin{theorem}[No finite-jet normality criterion]\label{thm:no-finite}
Fix $m\ge2$, a nonzero point $z_0$, and $K>0$.  There is an entire zero-free
function $F_K$, normalized by $F_K(0)=1$, whose logarithmic coefficients
through order $m$ all vanish, but $|F_K(z_0)|=e^K$.
Consequently no bound involving only a fixed logarithmic $m$-jet can imply
local boundedness on a neighborhood containing $z_0$.
\end{theorem}

\begin{proof}
Take
\[
 F_K(z)=\exp\!\left(K(z/z_0)^{m+1}\right).
\]
Its normalized logarithm begins at order $m+1$, while $F_K(z_0)=e^K$.
Letting $K\to\infty$ violates local boundedness.
\end{proof}

The theorem is deliberately formulated for zero-free entire functions, so
the obstruction persists inside the natural class of regularized determinant
quotients; it is not caused by inserting arbitrary poles or zeros.

\section{Two legitimate continuations}

There are two ways to strengthen the topology.  The direct route controls all
orders at once by a geometric envelope $|\tau_n|\le Mq^n$.  A second possible
route lets the audited order $m=m(\sigma)$ grow, but then it must also provide
a uniform tail estimate beyond $m(\sigma)$.  Merely increasing the computed
order without such a tail bound repeats Theorem~\ref{thm:no-finite} at a later
index.

Both routes must be combined with an anchor.  Convergence of the adaptive
jets to zero proves that the chosen quotient tends toward one in every fixed
audited coefficient; it does not prove that one is the intended deterministic
relative factor.

\section{Next wall}

RH-240 supplies the missing sufficient form: a geometric envelope
$|\tau_n(\sigma)|\le Mq^n$ for every $n\ge2$
\citep{WangRH240,Simon2005}.  The archived orders $2$--$12$
fit such an envelope with $q<1$, but order thirteen begins the unaudited tail.
In parallel, an anchoring theorem must show that the cloud has not absorbed
the regular deterministic numerator.

No Gate A closure or zeta-zero statement is made.

\bibliographystyle{plainnat}
\bibliography{references}
\end{document}
